% Figure for Relativity §C2.4 (push-forward and pull-back): a smooth map f : M -> N carries the point p, a curve c through p and its velocity V forward (f∘c, f_*V); a function g on N and a covector ω at f(p) are carried backward (g∘f, f^*ω). Bottom: the linear maps f_{*p} : T_pM -> T_{f(p)}N (forward) and f^*_p : T^*_{f(p)}N -> T^*_pM (backward). Schematic.
\begin{tikzpicture}[>={Stealth[length=2.2mm]}, font=\small, x=1cm, y=1cm]
  % ---------- manifold M ----------
  \draw[thick, black!75, fill=black!3] plot[smooth cycle, tension=0.8] coordinates {(-0.2,0.1) (1.3,-0.9) (3.1,-0.6) (3.6,0.9) (2.4,2.1) (0.5,1.8)};
  \node[font=\normalsize] at (0.3,2.05) {$M$};
  % curve c through p and its velocity V
  \draw[thick, figG] plot[smooth, tension=0.7] coordinates {(0.4,-0.1) (1.2,0.45) (1.8,0.65) (2.7,0.6) (3.2,0.35)};
  \fill (1.8,0.65) circle (1.6pt);
  \node[fill=black!3, inner sep=1pt] at (1.25,-0.25) {$p = c(0)$};
  \draw[black!55, thin] (1.4,-0.1) -- (1.76,0.58);
  \draw[->, very thick, figB] (1.8,0.65) -- (2.75,0.86) node[above, pos=0.85] {$V$};
  \node[figG, font=\footnotesize] at (0.55,0.45) {$c$};
  % pulled-back covector at p: level lines of g∘f (red dashed)
  \foreach \s in {-0.45,-0.15,0.15,0.45} \draw[figR, dashed, thin] ({1.8+\s-0.25},{0.65+0.75}) -- ({1.8+\s+0.25},{0.65-0.75});
  \node[figR, font=\footnotesize, align=center] at (1.5,1.65) {$f^{\ast}\omega$};
  % ---------- manifold N ----------
  \begin{scope}[xshift=6.6cm]
    \draw[thick, black!75, fill=black!3] plot[smooth cycle, tension=0.8] coordinates {(-0.4,0.3) (0.6,-1.0) (2.6,-0.9) (3.7,0.4) (3.0,2.0) (1.0,2.2)};
    \node[font=\normalsize] at (3.35,1.95) {$N$};
    \draw[thick, figG] plot[smooth, tension=0.7] coordinates {(0.2,-0.3) (0.9,0.2) (1.6,0.75) (2.2,1.35) (2.6,1.75)};
    \fill (1.6,0.75) circle (1.6pt);
    \node[fill=black!3, inner sep=1pt] at (1.85,-0.35) {$f(p)$};
    \draw[black!55, thin] (1.8,-0.2) -- (1.61,0.68);
    \draw[->, very thick, figB] (1.6,0.75) -- (2.28,1.52) node[right=7pt, pos=0.6] {$f_{\ast}V$};
    \node[figG, font=\footnotesize] at (0.25,0.15) {$f\circ c$};
    \foreach \s in {-0.4,-0.13,0.13,0.4} \draw[figR, dashed, thin] ({1.6+\s-0.55},{0.75+0.45}) -- ({1.6+\s+0.55},{0.75-0.45});
    \node[figR, font=\footnotesize] at (2.85,0.15) {$\omega$};
  \end{scope}
  % ---------- the real line and g ----------
  \draw[thick, black!75] (12.6,-0.9) -- (12.6,2.1);
  \node at (12.6,2.35) {$\mathbf R$};
  \fill (12.6,0.75) circle (1.4pt) node[right] {$g(f(p))$};
  % ---------- maps ----------
  \draw[->, thick] (3.7,1.5) to[bend left=18] node[above] {$f$} (6.3,1.5);
  \draw[->, thick] (10.4,1.0) to[bend left=12] node[above] {$g$} (12.45,1.0);
  \draw[->, thick, figR] (12.45,-0.35) to[bend left=14] node[below, pos=0.5] {$g\circ f = f^{\ast}g$} (3.3,-0.75);
  % ---------- linear maps at the points ----------
  \begin{scope}[yshift=-2.9cm]
    \node (TpM) at (1.8,0) {$T_pM$};
    \node (TfN) at (8.2,0) {$T_{f(p)}N$};
    \draw[->, thick, figB] (TpM) -- node[above] {$f_{\ast p}$: push-forward, $(f_{\ast}V)[g] = V[g\circ f]$} (TfN);
    \node (TsM) at (1.8,-1.1) {$T^{\ast}_pM$};
    \node (TsN) at (8.2,-1.1) {$T^{\ast}_{f(p)}N$};
    \draw[->, thick, figR] (TsN) -- node[below] {$f^{\ast}_p$: pull-back, $\langle f^{\ast}\omega, V\rangle = \langle\omega, f_{\ast}V\rangle$} (TsM);
    \node[font=\footnotesize, black!70, align=left, anchor=west] at (9.3,-0.55) {vectors move with $f$,\\ functions and covectors\\ move against it};
  \end{scope}
\end{tikzpicture}
